\subsection{Which Features for Which Goal}
\label{sec:choosing-features}
\slidesrc{Text features: 14--23; Blackboard: features}

\subsubsection{Content, FW, POS: when, why?}
\slidesrc{Text features: 14--22}

``Well\ldots\ it depends on what your goal is.''

\paragraph{Topic classification: content matters.}
If you need to classify a document into categories such as ``music'', ``biology'',
``food'' (three examples of BabelNet categories), you probably do not care about function
words and POS: what matters are the content words. Two texts from the slides:
\begin{itemize}[nosep]
  \item \emph{``Baked sea bass with lemon caper dressing. To make the dressing, mix the
        oil with the lemon zest and juice, capers, mustard, \ldots''} (a recipe: food);
  \item \emph{``These Russian basses sing so low you'll think the world is ending.
        \ldots\ the sound of a Russian oktavist, singing a full octave below a normal
        bass\ldots''} (music).
\end{itemize}
Both contain \emph{bass}, which in BabelNet has several senses: \emph{bass, basso} (an
adult male singer with the lowest voice), \emph{sea bass} (the lean flesh of a saltwater
fish), \emph{freshwater bass} (a fish), \emph{bass voice} (the lowest adult male singing
voice). The content words around it (\emph{baked}, \emph{lemon}, \emph{capers} vs.\
\emph{sing}, \emph{octave}) both pick the right sense (word sense disambiguation,
Section~\ref{sec:word-sense-disambiguation}) and the right topic. Counting pronouns or
articles would not help at all.

\paragraph{Author profiling: function words and POS matter.}
But if you are profiling the \emph{author} of a document\ldots

\begin{nbfact}{Argamon, Koppel, Pennebaker and Schler (2009)}
``However, the preponderance of evidence suggests that the most consistently effective
features over a wide variety of authorship problems are \textbf{function words} and
individual \textbf{parts-of-speech}.''

\smallskip\noindent
{\footnotesize S.\ Argamon, M.\ Koppel, J.\,W.\ Pennebaker, J.\ Schler,
\emph{Automatically profiling the author of an anonymous text}, Communications of the
ACM 52(2), 2009.}
\end{nbfact}

\begin{nbexample}{female vs.\ male}
The slides show two abstracts of linguistics papers, one by a female and one by a male
author, about similar topics (stylistic analysis, rhetorical devices). First the
\textbf{personal pronouns} are highlighted (\emph{my}, \emph{I}, \emph{he},
\emph{their}), then the preposition \textbf{of}. The point: the topic of the two texts
is the same, so content features cannot tell the authors apart, but the \emph{rates} of
such small words can. In the large studies by Argamon, Koppel and colleagues, texts by
women use more personal pronouns, texts by men more determiners, numbers and
prepositions such as \emph{of} (words that specify nouns).
\end{nbexample}

\subsubsection{So, what?}
\slidesrc{Text features: 23}

\begin{itemize}[nosep]
  \item There is a possibly endless amount of features. Many are standard and well
        studied, and it is also clear which NLP tasks they suit best; others (EmoGraph,
        for example) are less standard.
  \item Anyone can use existing features, already implemented by existing tools.
  \item The real challenge is to \textbf{invent new features}, or well-calibrated
        combinations of existing ones, and to \textbf{demonstrate} that they are useful
        for some given NLP task, on some existing benchmark.
\end{itemize}

\subsubsection{The blackboard summary}
\slidesrc{Blackboard: features}

In class the lecturer summarised which features are useful for which tasks in a table
on the blackboard. On the board the features are the \emph{columns}; here the table is
transposed (one feature per row) so that it fits the page. Nothing else is changed.

\begin{nbblackboard}{Which feature is useful for which task}{Blackboard: features}
\small
\renewcommand{\arraystretch}{1.25}
\begin{tabularx}{\linewidth}{@{}>{\bfseries\raggedright\arraybackslash}p{3.0cm}
  >{\raggedright\arraybackslash}X >{\raggedright\arraybackslash}p{4.0cm}@{}}
\toprule
\textnormal{\textbf{feature}} & \textbf{useful for} & \textbf{useless or not
  applicable for} \\
\midrule
Function words & author profiling: $\bullet$ gender of [the author] / pron[ouns];
  $\bullet$ one specific author & topic classification or filtering \\
TF-IDF & topic classification & not applicable without a corpus; WSD \\
POS & author profiling; genre profiling; readability & {\color{nbgray}\itshape(empty)} \\
Period / stylistic features & Bob or bot?; human author profiling; document genre
  profiling & {\color{nbgray}\itshape(empty)} \\
Emotions & sentiment analysis; affective computing; legal cases; fake news &
  {\color{nbgray}\itshape(empty)} \\
Length of words & {\color{nbgray}\itshape(empty)} & {\color{nbgray}\itshape(empty)} \\
Variability of words & readability & {\color{nbgray}\itshape(empty)} \\
Slang, flooding, jargon & {\color{nbgray}\itshape(empty)} &
  {\color{nbgray}\itshape(empty)} \\
\bottomrule
\end{tabularx}

\smallskip
{\color{nbgray}\scriptsize On the board a double arrow links the headings ``Period /
stylistic features'' and ``Length of words'' (they belong together). In the first cell
the text after ``gender of'' is partly unreadable; brackets mark what is completed from
context.}
\end{nbblackboard}

\paragraph{Reading the blackboard, row by row.}
\begin{itemize}
  \item \textbf{Function words}. Useful for \textbf{author profiling}: to guess the
        \emph{gender} of the author (through the use of pronouns, as in the female vs.\
        male example) or to recognise \emph{one specific author} (authorship
        attribution: everybody has a personal rate of \emph{the}, \emph{of},
        \emph{however}). Useless for \textbf{topic classification or filtering}: function
        words are the same in every topic, which is why they are removed as stop words.
  \item \textbf{TF-IDF}. Useful for \textbf{topic classification}: the words with high
        TF-IDF in a document are the ones that characterise its topic. \textbf{Not
        applicable without a corpus}: the idf needs a collection of documents $C$; with a
        single text we cannot compute it. Useless for \textbf{WSD}: TF-IDF weighs a word
        for a whole document, it says nothing about which sense one occurrence has.
  \item \textbf{POS}. Useful for \textbf{author profiling} (Argamon et al.),
        \textbf{genre profiling} (a scientific paper has more nouns, a novel more verbs
        and pronouns, a recipe many imperatives) and \textbf{readability} (many nouns and
        long noun phrases make a text harder).
  \item \textbf{Period / stylistic features}. ``Period'' here means the sentence
        (Italian \emph{periodo}): sentence length, punctuation, question sentences, the
        stylistic features of Table~\ref{tab:stylistic-features}. Useful to tell
        \textbf{a human (Bob) from a bot}, for \textbf{human author profiling}, and for
        \textbf{document genre profiling}.
  \item \textbf{Emotions}. Emotion words are useful for \textbf{sentiment analysis}
        (positive or negative opinion), \textbf{affective computing} (systems that
        recognise and react to human emotions), \textbf{legal cases} (e.g.\ the emotional
        load of testimonies or of threatening messages) and \textbf{fake news} (fake news
        often uses a strongly emotional, alarming language).
  \item \textbf{Length of words}. No task written on the board; it is linked by the
        arrow to the stylistic features (word length is one of them, features 5--7 of the
        stylistic table) and enters readability formulas.
  \item \textbf{Variability of words} (vocabulary richness, the number of unique words).
        Useful for \textbf{readability}: a text that uses many different and rare words
        is harder to read.
  \item \textbf{Slang, flooding, jargon}: no task written on the board; from the
        previous subsection they are idiosyncrasies of informal writers (useful for
        author profiling, e.g.\ age, and for bot detection).
\end{itemize}

\begin{nbintuition}[the rule behind the table]
Ask: \emph{does the task depend on what is said, or on who says it and how?}
\begin{itemize}[nosep]
  \item \textbf{What} is said (topic, filtering): use \emph{content} features (TF-IDF,
        semantic categories); remove function words.
  \item \textbf{Who} says it and \textbf{how} (author, gender, human vs.\ bot, genre,
        readability): use \emph{style} features (function words, POS, sentence and word
        length, vocabulary richness, idiosyncrasies).
  \item \textbf{How the author feels} (sentiment, fake news, affective computing): use
        \emph{emotion} features.
\end{itemize}
\end{nbintuition}

\begin{nbsummary}[choosing features]
\begin{itemize}[nosep]
  \item Topic classification: content words (TF-IDF); function words and POS useless.
  \item Author profiling: function words and POS are the most consistently effective
        features (Argamon et al., 2009); female vs.\ male: pronouns vs.\ noun specifiers.
  \item TF-IDF needs a corpus; useless for WSD.
  \item Stylistic (sentence-level) features: bot detection, author and genre profiling.
  \item Emotions: sentiment analysis, affective computing, legal cases, fake news.
  \item Variability of words, POS: readability.
  \item The challenge: invent new features and prove they help on a benchmark.
\end{itemize}
\end{nbsummary}
